% 05-modules.tex — Chapter 5: Module System
\chapter{Module System}
\label{chap:05-modules}
\index{module system}
\index{module loading}

\pnum
The Lain module system translates a flat dotted module name (e.g.\ \lain{foo.bar})
to a filesystem path (\texttt{foo/bar.ln}), reads and parses the file, and
recursively expands \lain{import} declarations in-place in the AST.  The entire
mechanism is implemented in \srcref{src/module.h}.

% -------------------------------------------------------------------------
\section{Module records}
\label{sec:mod-records}
\index{ModuleNode}
% -------------------------------------------------------------------------

\pnum
A \cname{ModuleNode} records every module that has been loaded during a compiler
run:

\begin{structdoc}{ModuleNode}{src/module.h:22}
\cname{name}        & \cname{char *}        & Canonical dotted name (\texttt{"foo.bar"}). \\
\cname{decls}       & \cname{DeclList *}    & Top-level declaration list for this module. \\
\cname{source\_text}& \cname{const char *}  & NUL-terminated source buffer (for diagnostic line display). \\
\cname{source\_file}& \cname{const char *}  & Filesystem path (for error messages). \\
\cname{next}        & \cname{ModuleNode *}  & Next node in the global linked list. \\
\end{structdoc}

\pnum
The global \cname{loaded\_modules} variable is the head of a singly-linked list
of all loaded modules.  Module records are allocated in \cname{ast\_arena}; the
name string is arena-copied so the record does not hold a pointer into a
stack-allocated buffer.

% -------------------------------------------------------------------------
\section{Module name to path conversion}
\label{sec:mod-path}
\index{module name!path conversion}
% -------------------------------------------------------------------------

\pnum
\cname{module\_name\_to\_path(mod, out, cap)} (\srcref{src/module.h:63})
translates a dotted module name to a filesystem path by replacing every
\texttt{.} with \texttt{/} and appending the \texttt{.ln} extension:

\begin{algorithm}
For each character \texttt{c} in \texttt{mod}:\\
\quad if \texttt{c == '.'}:\ write \texttt{/}\\
\quad else:\ write \texttt{c}\\
Append \texttt{".ln"}\,.
\end{algorithm}

\pnum
Paths are capped at 255 characters (the fixed \texttt{path[256]} buffer).
The compiler aborts if the module name is longer, which cannot occur for
reasonable module names.

% -------------------------------------------------------------------------
\section{load\_module algorithm}
\label{sec:mod-load}
\index{load\_module}
% -------------------------------------------------------------------------

\pnum
\cname{load\_module(file\_arena, ast\_arena, modname)}
(\srcref{src/module.h:78}) is the single entry point for module loading.
It is called recursively for each import dependency.

\begin{algorithm}
\textbf{load\_module}(\textit{file\_arena, ast\_arena, modname}):
\begin{enumerate}
\item If \texttt{modname} is already in \texttt{loaded\_modules}, return \texttt{NULL}.
  (Already loaded modules are silently skipped to prevent duplicate symbols.)
\item Translate \texttt{modname} to a filesystem path via
  \cname{module\_name\_to\_path}.
\item Read the file into \textit{file\_arena} using \cname{file\_read\_into\_arena}.
  Abort with an error if the file is missing.
\item Create a \cname{Lexer} pointing at the source buffer and a \cname{Parser}
  wrapping the lexer.
\item Call \cname{parse\_module(ast\_arena, \&parser)} to build a \cname{DeclList}.
\item Walk the top-level \cname{DeclList}.  For each \cname{DECL\_IMPORT} node:
  \begin{enumerate}[(a)]
  \item Recursively call \cname{load\_module} for the imported module name.
  \item If the recursive call returns a non-NULL \cname{DeclList}, splice it
    in-place into the current list: the import node is removed and the imported
    declarations are inserted at that position.
  \end{enumerate}
\item Call \cname{record\_module} to save the module record in the linked list.
\item Return the (now-spliced) \cname{DeclList}.
\end{enumerate}
\end{algorithm}

\pnum
The splice-in-place strategy means that after \cname{load\_module} returns, the
root \cname{DeclList} passed to \cname{sema\_resolve\_module} already contains
the declarations of all imported modules in a single flat list.  There is no
separate per-module processing in sema; all top-level symbols are globally visible
within the flat list.

\begin{rationale}
Splicing imports into a single flat list simplifies name resolution: there is no
module namespace or scope stack at the top level.  Every top-level symbol is
accessible from any function in the program.  This is sufficient for the current
Lain compiler which does not expose separate compilation or import qualified access.
\end{rationale}

% -------------------------------------------------------------------------
\section{Cycle prevention}
\label{sec:mod-cycles}
\index{module system!cycles}
% -------------------------------------------------------------------------

\pnum
Before loading a module, \cname{module\_already\_loaded()} performs a linear
scan of \cname{loaded\_modules}.  If the module is already present, \cname{NULL}
is returned immediately.  When the caller splices the result, a \texttt{NULL}
return means ``already loaded; nothing to splice'', so the import node is
silently dropped.

\begin{invariant}
Every distinct module name appears at most once in \cname{loaded\_modules}.
Mutually importing modules terminate because the second call to
\cname{load\_module("A")} from inside module B (which was loaded first from A)
finds A already registered and returns \texttt{NULL}.
\end{invariant}

% -------------------------------------------------------------------------
\section{Diagnostic integration}
\label{sec:mod-diagnostics}
% -------------------------------------------------------------------------

\pnum
\cname{record\_module} stores both the source text and the filesystem path in the
\cname{ModuleNode}.  Before running semantic analysis on a module,
\cname{sema\_resolve\_module()} calls \cname{find\_module(module\_path)} and
copies the pointers into the globals \cname{sema\_source\_text} and
\cname{sema\_source\_file} (\srcref{src/sema.h:841}).  These are then available
to \cname{diagnostic\_show\_line()} when any semantic error is reported.

\pnum
The \texttt{path} variable used as \cname{source\_file} is the local
\texttt{char path[256]} buffer inside \cname{load\_module}.  Its address is
stored in the module record, which means the module record is invalidated if
the stack frame unwinds.  This is safe only because \cname{record\_module} is
called at the end of \cname{load\_module} before the function returns, and the
caller uses the pointer while the frame is still live, or — after the refactor
that allocates records in \cname{ast\_arena} — the pointer points to a fixed
path variable whose lifetime is bounded by the module loading phase.

\begin{note}
A future improvement would copy \texttt{path} into \cname{ast\_arena} before
storing it, eliminating the dependency on stack lifetime.
\end{note}
